%!TEX program = xelatex
% ################# recalculate final print size ############ 
\documentclass[12pt]{article}
% full height = 297mm
% full width = 210 + 210 + spine width = 444.5mm
% spine width = 24.5mm
% https://www.papersizes.org/a-paper-sizes.htm?utm_expid=76296969-1675.RJoMUqt0TyO9xVbJnRsT2w.0&utm_referrer=https%3A%2F%2Fwww.google.fr%2F
% front/back cover width = 8in
\usepackage[paperheight=297mm,paperwidth=444.5mm,margin=0mm]{geometry}
\setlength{\parindent}{0pt}
\usepackage[dvipsnames,prologue,table]{pstricks}
\usepackage{graphicx}
\usepackage[T1]{fontenc}
\usepackage{lmodern}
\usepackage{pst-text}% http://ctan.org/pkg/pst-text
\usepackage[scaled]{helvet}
\renewcommand*\familydefault{\sfdefault}
% \usepackage{uarial}
% \renewcommand{\familydefault}{\sfdefault}

%R147 G191 B235
\definecolor{fcup}{rgb}{0.57,0.75,0.92}

\usepackage{lipsum}

\usepackage{environ}% http://ctan.org/pkg/environ
\newdimen\fontdim
\newdimen\upperfontdim
\newdimen\lowerfontdim
\newif\ifmoreiterations
\fontdim12pt

\makeatletter
\NewEnviron{fitbox}[2]{% \begin{fitbox}{<width>}{<height>} stuff \end{fitbox}
  \def\buildbox{%
    \setbox0\vbox{\hbox{\minipage{#1}%
      \fontsize{\fontdim}{1.2\fontdim}%
      \selectfont%
      \stuff%
    \endminipage}}%
    \dimen@\ht0
    \advance\dimen@\dp0
  }
  \def\stuff{\BODY}% Store environment body
  \buildbox
  % Compute upper and lower bounds
  \ifdim\dimen@>#2
    \loop
      \fontdim.5\fontdim % Reduce font size by half
      \buildbox
    \ifdim\dimen@>#2 \repeat
    \lowerfontdim\fontdim
    \upperfontdim2\fontdim
    \fontdim1.5\fontdim
  \else
    \loop
      \fontdim2\fontdim % Double font size
      \buildbox
    \ifdim\dimen@<#2 \repeat
    \upperfontdim\fontdim
    \lowerfontdim.5\fontdim
    \fontdim.75\fontdim
  \fi
  % Now try to find the optimum size
  \loop
    %\message{Bounds: \the\lowerfontdim\space
    %         \the\fontdim\space \the\upperfontdim^^J}
    \buildbox
    \ifdim\dimen@>#2
      \moreiterationstrue
      \upperfontdim\fontdim
      \advance\fontdim\lowerfontdim
      \fontdim.5\fontdim
    \else
      \advance\dimen@-#2
      \ifdim\dimen@<10pt
        \lowerfontdim\fontdim
        \advance\fontdim\upperfontdim
        \fontdim.5\fontdim
        \dimen@\upperfontdim
        \advance\dimen@-\lowerfontdim
        \ifdim\dimen@<.2pt
          \moreiterationsfalse
        \else
          \moreiterationstrue
        \fi
      \else
        \moreiterationsfalse
      \fi
    \fi
  \ifmoreiterations \repeat
  \box0% Typeset content
}
\makeatother


%% This code is to make empty letters for the 'PhD's in the corners
\usepackage{stackengine,xcolor}
\newcommand\shadowfy[1]{\expandafter\shadowfypars#1\par\relax\relax}
\long\def\shadowfypars#1\par#2\relax{%
  \ifx#1\relax\else
    \shadowfywords#1 \relax\relax%
  \fi%
  \ifx\relax#2\else\par\shadowfypars#2\relax\fi%
}
\def\shadowfywords#1 #2\relax{%
  \ifx#1\relax\else
    \shadowfyletters#1\relax\relax%
  \fi%
  \ifx\relax#2\else\ \shadowfywords#2\relax\fi%
}
\def\shadowfyletters#1#2\relax{%
  \shadow{#1}%
  \ifx\relax#2\else\shadowfyletters#2\relax\fi}

\newlength\shadowHoffset
\newlength\shadowVoffset
\setlength\shadowHoffset{.2pt}
\setlength\shadowVoffset{.1pt}
\def\primarycolor{white}
\def\secondarycolor{black}

\def\shadow#1{\setstackgap{L}{0pt}\def\stacktype{L}%
  \def\useanchorwidth{T}% CAN BE COMMENTEDD FOR MORE INTERLETTER SPACE.
  \Longstack{%
  \raisebox{0pt}{\textcolor{\primarycolor}{#1}} 
  \kern.7\shadowHoffset\raisebox{.7\shadowVoffset}{\textcolor{\secondarycolor}{#1}}
  \kern-.7\shadowHoffset\raisebox{.7\shadowVoffset}{\textcolor{\secondarycolor}{#1}}
  \kern\shadowHoffset\raisebox{0pt}{\textcolor{\secondarycolor}{#1}}
  \kern-\shadowHoffset\raisebox{0pt}{\textcolor{\secondarycolor}{#1}}
  \kern.7\shadowHoffset\raisebox{-.7\shadowVoffset}{\textcolor{\secondarycolor}{#1}}
  \kern-.7\shadowHoffset\raisebox{-.7\shadowVoffset}{\textcolor{\secondarycolor}{#1}}
  \kern0pt\raisebox{\shadowVoffset}{\textcolor{\secondarycolor}{#1}}
  \kern0pt\raisebox{-\shadowVoffset}{\textcolor{\secondarycolor}{#1}}%
}}



\begin{document}
\thispagestyle{empty}
\psset{unit=1in}
\begin{pspicture}(0,0)(\paperwidth,11.69in)
% \psgrid


\newsavebox\IBoxBig
\sbox\IBoxBig{\includegraphics[height=1in]{fcup_clipped.eps}}
\newsavebox\IBoxSmall
\sbox\IBoxSmall{\includegraphics[height=0.3in]{uporto_inverted.eps}}
\newsavebox\FCBox
\sbox\FCBox{\includegraphics[height=0.3in]{fc.eps}}


\rput[Bl](0,5.345){\usebox\IBoxBig}
\newlength{\xposlogo}
\setlength{\xposlogo}{\paperwidth - 6cm - 215pt}
\rput[Bl](\xposlogo, 10){\usebox\IBoxBig}


\setlength{\fboxsep}{100pt}
\setlength{\fboxrule}{2pt}
\newlength{\xpostitle}
\setlength{\xpostitle}{\paperwidth - 12cm - 215pt}
\rput[tl](\xpostitle,7){
% The fontsize of the title adjusts to the 12cm x 12cm box
\begin{fitbox}{12cm}{12cm}
%%% CHANGE THIS %%%
% feel free (and encouraged!) to do line breaks where you want
\textbf{Obstacle \\ Avoidance \\ Framework \\ based on \\ Reach Sets}
\end{fitbox}}


%% depending on how big the title is,
%% you can shift this text up or down by changing the \vpos value
\newlength{\vpos}
\setlength{\vpos}{1.8in}
\setlength{\xpostitle}{\xpostitle + 0.2cm} %% no idea what the 0.2cm is...

\rput[Bl](\xpostitle,\vpos){{\fontsize{18pt}{1em}\selectfont %
  %%% CHANGE THIS %%%
  Alojz Gomola
}}

\setlength{\vpos}{\vpos - .4in}%
\rput[Bl](\xpostitle,\vpos){{\fontsize{14pt}{1em}\selectfont Tese de Doutoramento apresentada à}}
\setlength{\vpos}{\vpos - .2in}%
\rput[Bl](\xpostitle,\vpos){{\fontsize{14pt}{1em}\selectfont Faculdade de Ciências da Universidade do Porto}}
\setlength{\vpos}{\vpos - .2in}%
\rput[Bl](\xpostitle,\vpos){{\fontsize{14pt}{1em}\selectfont Matemática Aplicada}}
\setlength{\vpos}{\vpos - .4in}%
\rput[Bl](\xpostitle,\vpos){{\fontsize{18pt}{1em}\selectfont 2019}}



\newlength{\posmiddle}
\setlength{\posmiddle}{\paperwidth / \real{2.0}}
%put the text on the spine (note the rotation over 270 degrees)
\rput[l]{270}(\posmiddle,11.68){{\fontsize{8.45mm}{1em}\selectfont PhD}}

\rput[l]{270}(\posmiddle, 7.5){{%
    \fontsize{11pt}{1em}\selectfont %
      %%% CHANGE THIS %%%
      \textbf{Obstacle Avoidance Framework based on Reach Sets}
    }}


\rput[ml]{270}(\posmiddle, 3.5){{%
    \fontsize{10pt}{1em}\selectfont %
      %%% CHANGE THIS %%%
      \textbf{Alojz Gomola}
    }}


%put the logos in the spine
\rput[c]{270}(\posmiddle,8.5){{\usebox\IBoxSmall}}
\rput[c](\posmiddle,.5){{\usebox\FCBox}}


% the spine length is 1.05in
% but we make this blue box 1.5in wide
% 17.59/2 + 1.5/2 = 9.545
\newlength{\xposone}
\setlength{\xposone}{\paperwidth / \real{2.0} - 0.75in}
\newlength{\xpostwo}
\setlength{\xpostwo}{\paperwidth / \real{2.0} + 0.75in}

\psframe*[linewidth=0.01,linecolor=fcup](\xposone,10.4)(\xpostwo,10.8)

\rput[c](\posmiddle, 10.5){{%
    \fontsize{10pt}{1em}\selectfont \textcolor{white}{\textbf{CICLO}}}}
\rput[c](\posmiddle, 10.65){{%
    \fontsize{10pt}{1em}\selectfont \textcolor{white}{3$^{\circ}$}}}

\rput[c](\posmiddle,10.2){{\fontsize{10pt}{1em}\selectfont FCUP}}
\rput[c](\posmiddle,10.0){{\fontsize{10pt}{1em}\selectfont 2019}}



%%% options for the outlined text in the borders
\setlength\shadowHoffset{5pt}
\setlength\shadowVoffset{2pt}
\def\primarycolor{white}
\def\secondarycolor{fcup}

% top right corner
\rput[tl]{270}(\paperwidth,11.9){%
  {\fontsize{215pt}{1em}\selectfont %
    \textcolor{fcup}{\shadowfy{PhD}}
  }
}
% bottom right corner
\rput[lb]{90}(\paperwidth,-0.2){%
  {\fontsize{215pt}{1em}\selectfont %
    \textcolor{fcup}{\shadowfy{Ph}D}
  }
}
% bottom left corner
\rput[lb]{-90}(0, 5.3){%
  {\fontsize{215pt}{1em}\selectfont %
    \textcolor{fcup}{\reflectbox{\shadowfy{Ph}D}}
  }
}



\end{pspicture}
\end{document}
